\documentclass[10pt,a4paper]{amsart}
\usepackage[margin=19mm]{geometry}
\usepackage{amsmath,amssymb,mathtools}
\usepackage[scr=rsfs,cal=euler]{mathalfa}
\usepackage{xfrac}
\usepackage[unicode,hidelinks,bookmarks=false]{hyperref}
\newcommand{\ie}{i.e.\ }
\newcommand{\cf}{cf.\ }
\newcommand{\resp}{respectively\ }
\newcommand{\Subsection}{\subsection}
\providecommand{\nobreakdash}{\nobreak\hbox{-}}
\setlength{\emergencystretch}{2em}
\multlinegap0pt
\usepackage{amssymb}
\usepackage{dsfont}
\newtheorem{lem}{Lemma}
\title{On semi-openness of fiber-onto extensions of minimal semiflows and quasi-separable maps}
\author{Xiongping Dai, Li Feng, Congying Lv, Yuxuan Xie}
\date{Source: arXiv:2601.03380v2 (CC BY 4.0)}
\begin{document}
\maketitle

\noindent\emph{Source and licence.} On semi-openness of fiber-onto extensions of minimal semiflows and quasi-separable maps. Version arXiv:2601.03380v2 (CC BY 4.0), distributed by arXiv under the Creative Commons Attribution 4.0 International licence, http://creativecommons.org/licenses/by/4.0/, as recorded in the arXiv licence metadata for this exact version and shown on the abstract page https://arxiv.org/abs/2601.03380v2. Full licence notice: \texttt{licenses/arxiv-2601.03380-CC-BY-4.0.txt}.\par\smallskip

\noindent\textit{Editorial scope.} Counted unit: the article's lem 4.4 (source0.tex lines 812--814). The article's complete original proof of it is delivered with it (source0.tex lines 816--827, one \texttt{proof} environment of 12 lines). No mathematics is written, repaired or re-derived: the statement and the proof are the article's own lines, in the article's own notation, and no retained line is substituted or re-flowed.  No further source-local context is needed: every cross-reference of every retained line resolves inside the two delivered spans.  Nothing was omitted from any retained span, so the package carries no omission record.  No retained line contains a citation, so the wrapper carries no bibliography and no bibliography entry.  No retained line contains a figure environment, an image-inclusion command or drawing code, so no image is delivered and the package declares no asset dependency.  Editorial formatting only, in addition: an AMS \texttt{amsart} preamble that restates the article's own declarations and macros used by the retained lines, namely the environments \texttt{lem} on separate counters, together with the standard packages those lines need.  The retained environments print in this document as Lemma 1 in order of appearance, whereas the article prints the same items with the numbering of its own full document.  The complete native source of the article as distributed in the arXiv e-print for this exact version is bundled unchanged as \texttt{sources/192221/source0.tex}.
\begin{lem}\label{4.4}
If $\{X_i\,|\,i\in\Lambda\}$ is an inverse system of compact Hausdorff spaces, then we have that $\mathcal{M}^1\left(\underleftarrow{\lim}_{i\in\Lambda}\{X_i\}\right)\cong\underleftarrow{\lim}_{i\in\Lambda}\left\{\mathcal{M}^1(X_i)\right\}$.
\end{lem}

\begin{proof}
Write $X_i\xleftarrow{\pi_{i,j}}X_j$ and $\pi_{i,i}=\textsl{id}_{X_i}$, for $i<j$ in $\Lambda$, for the link maps of the inverse system $\{X_i\,|\,i\in\Lambda\}$. Let $X=\underleftarrow{\lim}\{X_i\}$. Let $p_i\colon X\rightarrow X_i$ and $p_{i*}\colon\mathcal{M}^1(X)\rightarrow\mathcal{M}^1(X_i)$ be the canonical maps. As $p_i=\pi_{i,j}\circ p_j$ for $i<j$ in $\Lambda$, it follows that $p_{i*}=\pi_{i,j*}\circ p_{j*}$ so that $\{\mathcal{M}^1(X_i)\,|\,i\in\Lambda\}$ is an inverse system.
Moreover,
$\pi\colon\mathcal{M}^1(X)\rightarrow\underleftarrow{\lim}\{\mathcal{M}^1(X_i)\}$, $\mu\mapsto(p_{i*}(\mu))_{i\in\Lambda}$,
is a continuous onto map. In fact, for all $(\mu_i)_{i\in\Lambda}\in\underleftarrow{\lim}\{\mathcal{M}^1(X_i)\}$, by the finite-intersection property of compact space we may take a probability $\mu\in\bigcap_{i\in\Lambda}{p_i}_*^{-1}(\mu_i)\subseteq\mathcal{M}^1(X)$; then $\pi(\mu)=(\mu_i)_{i\in\Lambda}$.
Now let $\lambda, \mu\in\mathcal{M}^1(X)$ such that $p_{i*}(\lambda)=:\lambda_i=\mu_i:=p_{i*}(\mu)$ for all $i\in\Lambda$. We need prove that $\lambda=\mu$. By regularity of $\lambda$ and $\mu$, it is enough to show that $\lambda(K)=\mu(K)$ for all $K\in2^X$. Let $K\in2^X$ and $\varepsilon>0$. Then we can take a set $U\in\mathscr{O}(X)$ such that $\lambda(U\setminus K)+\mu(U\setminus K)<\varepsilon$ and $K\subset U$. In addition, as $i\in\Lambda$ sufficiently big, we can choose finitely many open sets, say $V_1,\dotsc,V_\ell$ in $X_i$ such that $K\subseteq p_i^{-1}(V_1)\cup\dotsm\cup p_i^{-1}(V_\ell)\subseteq U$. So by the inclusion-exclusion formula of probability or by the equality $p_i^{-1}(V_1)\cup\dotsm\cup p_i^{-1}(V_\ell)=p_i^{-1}(V_1\cup\dotsm\cup V_\ell)$, it follows that
\begin{equation*}
\lambda(K)\le\lambda(p_i^{-1}(V_1)\cup\dotsm\cup p_i^{-1}(V_\ell))
=\mu(p_i^{-1}(V_1)\cup\dotsm\cup p_i^{-1}(V_\ell))\le\mu(K)+\varepsilon.
\end{equation*}
Thus, $\lambda(K)\le\mu(K)$; and analogously, $\mu(K)\le\lambda(K)$. Then $\pi$ is 1-1 onto and $\pi$ is a homeomorphism. The proof is complete.
\end{proof}

\end{document}
